\section{Tabelas, matemática e referências}

\begin{frame}[fragile]{Tabela básica}
\begin{columns}[T]
\begin{column}{0.55\textwidth}
\begin{lstlisting}
\begin{tabular}{lcc}
  \toprule
  Modelo & Treino & Validacao \\
  \midrule
  ResNet    & 0,94 & 0,90 \\
  MobileNet & 0,91 & 0,89 \\
  \bottomrule
\end{tabular}
\end{lstlisting}
\end{column}
\begin{column}{0.41\textwidth}
\centering
\begin{tabular}{lcc}
  \toprule
  Modelo & Treino & Val. \\
  \midrule
  ResNet & 0,94 & 0,90 \\
  MobileNet & 0,91 & 0,89 \\
  \bottomrule
\end{tabular}
\end{column}
\end{columns}

Em \verb|{lcc}|, cada letra cria uma coluna: esquerda (\texttt{l}) ou centralizada (\texttt{c}). O caractere \verb|&| separa células e \verb|\\| encerra a linha.
\end{frame}

\begin{frame}[fragile]{Expressões matemáticas}
\begin{columns}[T]
\begin{column}{0.52\textwidth}
\begin{lstlisting}
Texto com $3 \times 3$ na linha.

\[
  acuracia = \frac{acertos}{total}
\]

\begin{equation}
  y = f(Wx + b)
\end{equation}
\end{lstlisting}
\end{column}
\begin{column}{0.44\textwidth}
Texto com $3 \times 3$ na linha.

\[
  \mathrm{acurácia} = \frac{\mathrm{acertos}}{\mathrm{total}}
\]

\begin{equation}
  y = f(Wx + b)
\end{equation}
\end{column}
\end{columns}
\end{frame}

\begin{frame}[fragile]{Citações no texto}
\begin{lstlisting}
Segundo \textcite{LeCun2015}, redes profundas
aprendem representacoes em varias camadas.

O tema tambem pode aparecer entre parenteses
\cite{Goodfellow-et-al-2016}.

\nocite{8308186}
\end{lstlisting}

\begin{itemize}
    \item \verb|\textcite| inclui o autor na frase.
    \item \verb|\cite| produz uma citação entre parênteses.
    \item \verb|\nocite| envia uma fonte à bibliografia sem citá-la no texto.
\end{itemize}
\end{frame}

\begin{frame}[fragile]{Arquivo bibliográfico}
No preâmbulo:
\begin{lstlisting}
\usepackage[style=abnt]{biblatex}
\addbibresource{ref.bib}
\end{lstlisting}

No ponto onde a lista deve aparecer:
\begin{lstlisting}
\begin{frame}[allowframebreaks]{Referencias}
  \printbibliography
\end{frame} % encerra o slide
\end{lstlisting}

\verb|allowframebreaks| autoriza o Beamer a dividir uma bibliografia longa em vários slides.
\end{frame}
